\paragraph{The analytic positivity theorem and its arithmetic hypothesis.}
K\"unnemann, \emph{Some remarks on the arithmetic Hodge index
conjecture}, Compositio Mathematica 99 (1995), 109--128,
\url{https://www.numdam.org/item/CM_1995__99_2_109_0.pdf},
separates two statements that must not be conflated in the present
problem. Theorem~1.2 (p.~111) gives positivity for the specified
analytic pairing on a quotient of smooth differential forms. In
\S5, Conjecture~$D_p$, equation~(32), p.~124, asks for the signed
positivity of a canonical cubical height of the correspondence
$(\operatorname{id}\times\varphi)_*(Z\times Z)$ for a nonzero
homologically trivial cycle. Theorem~5.1 on the same page assumes
$D_p$, together with its stated $B$ and $C$ hypotheses, to obtain
the arithmetic Hodge-index conclusion. The Fourier calculation on
p.~125, notably equation~(36), identifies the grade-one contribution
with that height condition; it does not prove $D_p$ from Theorem~1.2.

Consequently combining the accepted finite-field algebraic Hodge
positivity with this analytic theorem leaves precisely an arithmetic
height comparison. AR-10 provides its canonical adelic line and
AR-18--AR-26 sharpen the comparison, but neither the analytic positive
form nor the separate nonnegative fibre forms supplies the missing
signed global degree. This is a source-scope check, not a new
counterexample to the arithmetic conjecture. The original printed
pp.~124--125 were visually inspected in this run.
